%--------------------------------------------------------------------------------
%
% 
%--------------------------------------------------------------------------------
\unnumberedChapter{A Layout Control System}

At a higher level, the layout control system consists of components and a communication scheme. This chapter will define the key concepts of a layout system. At the heart of the layout control system is a common communication bus to which all modules connect. The others key elements are node, events, ports and attributes. Let's define these items first and then talk about how they interact. The following figure depicts the high level view of a layout control system.

\begin{center}
    \begin{tikzpicture}[scale=0.9, transform shape]
        
        % \draw[help lines, gray!50, dashed] (0,0) grid(16,9);

        % Thick horizontal line (LCS bus) with text
        \draw[line width=1.5mm, <->, line cap=round, draw=gray, name path=lcsline] 
            (0,6) -- (16,6)
            node[pos= 0.9, tsLargeBold, below] {LCS bus};

        % Define nodes
        \node[  tsRoundedRectangle, 
                minimum width=3cm,
                minimum height=1.5cm,
                text width=3cm,
                text centered,
                fill=red!20] (gw) at (3,8) {Gateways};

        \node[  tsRoundedRectangle, 
                minimum width=3cm,
                minimum height=1.5cm,
                text width=3cm,
                text centered,
                fill=red!20] (bs) at (7,8) {Base\\Station};

        \node[  tsRoundedRectangle, 
                minimum width=3cm,
                minimum height=1.5cm,
                text width=3cm,
                text centered,
                fill=red!20] (ch) at (13,8) {Cab\\Handheld};

        % Connecting arrows to LCS bus
        \path[name path=p1] (gw.south) -- (3,6); % Define arrow path
        \path[name intersections={of=lcsline and p1, by=intpoint1}];
        \coordinate (adjustedPoint1) at ($(intpoint1) + (0,0.75mm)$);
        \draw[<->, ultra thick, line cap=round] (gw.south) -- (adjustedPoint1);

        % Arrow from Base Station
        \path[name path=p2] (bs.south) -- (7,6); % Define arrow path
        \path[name intersections={of=lcsline and p2, by=intpoint2}]; % Find intersection
        \coordinate (adjustedPoint2) at ($(intpoint2) + (0,0.75mm)$);
        \draw[<->, ultra thick, line cap=round] (bs.south) -- (adjustedPoint2); % Draw to intersection

        % Arrow from Cab Handheld
        \path[name path=p3] (ch.south) -- (13,6);
        \path[name intersections={of=lcsline and p3, by=intpoint3}];
        \coordinate (adjustedPoint3) at ($(intpoint3) + (0,0.75mm)$);
        \draw[<->, ultra thick, line cap=round] (ch.south) -- (adjustedPoint3);
        
        \node[ 	tsRoundedRectangle,
    			minimum width=12cm,
    			minimum height=4cm,
    			fill=yellow!20
				] (a) at (8,2.5) {};

		\node[	tsLargeBold,
				anchor=south,
    			align=center,
    			yshift=3mm,
			 ] at (a.south) {Layout};
        
        % Other nodes
        \node[  tsRoundedRectangle,  
                minimum width=3cm,
                minimum height=1.5cm,
                text width=3cm,
                text centered,
                fill=yellow!20] (bc) at (4,3) {Block\\Controllers};

        \node[  tsRoundedRectangle, 
                minimum width=3cm,
                minimum height=1.5cm,
                text centered,
                text width=3cm, 
                fill=yellow!20] (s) at (8,3) {Sensors};

        \node[  tsRoundedRectangle, 
                minimum width=3cm,
                minimum height=1.5cm,
                text width=3cm,
                text centered, 
                fill=yellow!20] (a) at (12,3) {Actors};

        % Vertical arrows to LCS bus \path[name path=p3] (bc.north) -- (4,8);
        \path[name path=p4] (bc.north) -- (4,8);
        \path[name intersections={of=lcsline and p4, by=intpoint4}];
        \coordinate (adjustedPoint4) at ($(intpoint4) - (0,0.75mm)$);
        \draw[<->, ultra thick, line cap=round] (bc.north) -- (adjustedPoint4);

        \path[name path=p5] (s.north) -- (8,8);
        \path[name intersections={of=lcsline and p5, by=intpoint5}];
        \coordinate (adjustedPoint5) at ($(intpoint5) - (0,0.75mm)$);
        \draw[<->, ultra thick, line cap=round] (s.north) -- (adjustedPoint5);

        \path[name path=p6] (a.north) -- (12,8);
        \path[name intersections={of=lcsline and p6, by=intpoint6}];
        \coordinate (adjustedPoint6) at ($(intpoint6) - (0,0.75mm)$);
        \draw[<->, ultra thick, line cap=round] (a.north) -- (adjustedPoint6);
        
        
                
    \end{tikzpicture}
\end{center}

\unnumberedSection{LCS Bus}

The layout control system is organized around a common communication bus. This bus is the backbone of the system: every hardware module connects to it and exchanges messages with the other modules. The current implementation uses the industry-standard CAN bus.

CAN messages are broadcast on the bus, so every connected module can receive them. Under the classic CAN standard, a message payload is limited to eight bytes; consequently, eight bytes is also the maximum message size defined for the LCS bus. In a practical installation, available bandwidth limits a single CAN bus to roughly 110 modules. This is sufficient for most layouts. Very large layouts could use an alternative bus technology or several CAN segments connected by routers.

The software architecture should nevertheless be able to manage hundreds of modules. Although the underlying transport technology may eventually change, the LCS message format, maximum message size, and broadcast communication model are part of the system design and remain stable.

\unnumberedSection{LCS Hardware}

Every physical device connected to the LCS bus is a \textbf{hardware module}. A module normally combines a microcontroller, a CAN-bus interface, non-volatile storage, and purpose-built electronics. For example, a module that controls turnouts and signals might consist of a Raspberry Pi Pico controller, a CAN interface, and digital-output drivers. Base stations, cab handhelds, gateways, and block controllers are further examples.

Hardware modules are normally installed close to the equipment they serve and are therefore distributed throughout the layout. Some may be in locations that are difficult to reach. Configuration and operation must consequently be possible through LCS messages alone, although local controls on a module may still be useful and are not excluded.

Examples of hardware modules include:

\begin{smallGapItemize}
\item \textbf{Main controller}. A general purpose controller without any special function other than the CAN bus interface and non-volatile storage.
\item \textbf{Base Station}. A controller with a power unit to generate DCC signal and manage locomotive sessions.
\item \textbf{Block Controller}. A controller with a power unit to drive DCC and PWM signals to the tracks.

\end{smallGapItemize}

Each hardware module has two conceptual parts: the \textbf{module main controller} and the \textbf{module extension}. The main controller contains the controller chip, non-volatile memory for retaining data across power cycles, the CAN-bus interface, and interfaces to the node-specific electronics. The node-specific electronics are provided by the module extension.

These two parts may be implemented together on a single printed-circuit board or as separate boards joined by a standardized extension connector. The main controller provides the LCS-bus interface and power for itself and its extensions. It also provides a local I\textsuperscript{2}C bus, which is the primary means of communication with the extension boards. The standardized connector carries the required power and communication signals and allows up to four extension boards to be connected to one main controller. The hardware chapters describe the physical board layouts and design options in more detail.

This separation provides a practical balance between common infrastructure and specialized functions. A common main controller can be paired with different extensions, allowing a module to be tailored to the sensors or actors required at a particular location on the layout.

\unnumberedSection{LCS Hardware Extensions}

Sensors and actors are the eyes, ears, and hands of a layout control system. Their possible applications are numerous, and the list of desirable functions is never truly complete. Sensors are primarily event producers: examples include block-occupancy detection, power-overload detection, and a push button on a layout control panel. Actors are their counterparts. They are the hardware that acts on the layout---for example, by operating turnouts, signals, relays, or other devices.

The variety of sensors and actors makes reusability essential. The main-controller-and-extension architecture permits functional building blocks to be combined in different ways without redesigning the common controller infrastructure. The following chapters describe common sensor and actor nodes and show how they can be assembled from these basic building blocks.

Examples of extension hardware include:

\begin{smallGapItemize}
\item \textbf{Occupancy Detector}. This extension is a companion to the block-controller board. It provides four track channels that correspond to the block-controller outputs, with a set of detectors on each channel.
\item \textbf{Turnout}. This extension controls a set of turnouts and provides frog polarization. Like the occupancy-detector extension, it is intended for use with the block-controller board.
\item \textbf{Signal}. This general-purpose extension drives LED-equipped signals. It meets their power requirements and supports functions such as dimming and blinking of individual lamps.
\item \textbf{Servo}. This universal extension can be connected to any of the controller boards presented in this book. It provides a set of general-purpose servo outputs.
\item \textbf{GPIO}. This universal digital input/output extension is typically used for switches and push buttons as inputs, and for LEDs as outputs. The number and type of output loads must be chosen within the drive capability of the individual output pins.
\item \textbf{Relays}. This general-purpose extension drives a set of external relays. It can support relays with different coil-voltage requirements.
\item \textbf{Cab Control}. Cab handhelds and stationary cab-control devices can also be implemented as extensions. For example, such an extension can be combined with a base station to create a complete control station for a smaller layout.
\item \textbf{Prototyping}. A prototyping extension provides the extension-address-resolution logic and space for a custom hardware design. It is useful for quickly developing or testing a new function.
\end{smallGapItemize}

\unnumberedSection{LCS Firmware}

A hardware module is the physical implementation; its firmware is the software that runs on it. The firmware is organized in layers so that module-specific functions can be developed independently of the details of a particular controller family.

The lowest software layer above the hardware is the \textbf{controller-dependent code}. It isolates the upper layers from controller-specific details and presents a higher-level interface to the bare-metal functions. Direct hardware access is still possible when it is necessary, but it should be the exception rather than the default.

Above this layer sits the \textbf{LCS Runtime Library}. The runtime provides the common services of the layout control system, including hardware-module discovery, attribute management, message passing, and event processing. It also manages non-volatile storage and handles the regular processing of LCS messages. Messages that require module-specific action are delivered to the upper firmware layer; routine protocol work and periodic housekeeping are performed by the runtime. This frees the module developer from much of the repetitive infrastructure code.

At the top are the module-specific functions. A base station, for example, implements the DCC subsystem and locomotive management. Other module types provide their own specialized functions while using the services of the runtime library and controller-dependent code. Where justified, they may also access the underlying hardware directly.

\textbf{Callbacks} are a central part of this model. A typical firmware initializes the runtime, which in turn initializes the lower layers. It then registers callback routines for the runtime events relevant to the module. Finally, the firmware starts the runtime. From that point, the runtime receives messages, performs periodic tasks, and invokes the registered callbacks whenever the module must respond.
